\documentclass[11pt]{article}
\usepackage[margin=1in]{geometry}
\usepackage{amsmath,amssymb,amsthm,hyperref}
\newtheorem{theorem}{Theorem}
\title{A strictly positive definite counterexample to coefficient-minor stability}
\author{Formal verification of a posted counterexample by Shivam Patel}
\date{October 26, 2026}
\begin{document}\maketitle
\begin{abstract}
We verify an exact counterexample to Fisk's three-matrix coefficient-minor stability conjecture, recorded as MathDB 380287. Three strictly positive definite integer matrices of order three yield a coefficient minor having a positive real zero. Lean certifies the positivity hypotheses, the determinant and coefficient extraction, and the obstruction to Hurwitz stability. This is verification of an existing posted counterexample, with no discovery or priority claim.
\end{abstract}
Fisk's Conjecture 21.90 asks whether, for positive definite matrices $D_1,D_2,D_3$, every minor of the coefficient matrix
\[
 M_{ij}(x)=[y^iz^j]\det(I+xD_1+yD_2+zD_3),\qquad i,j\geq0,
\]
is stable. At the start of Chapter 21 the source defines stability by requiring all roots to lie in the closed left half-plane. Equivalently the polynomial is nonzero at every complex number with strictly positive real part.
\begin{theorem}
Let
\[
D_1=\begin{pmatrix}1&0&0\\0&2&0\\0&0&1\end{pmatrix},\quad
D_2=\begin{pmatrix}13&9&7\\9&7&5\\7&5&4\end{pmatrix},\quad
D_3=\begin{pmatrix}5&7&8\\7&11&12\\8&12&14\end{pmatrix}.
\]
All three matrices are strictly positive definite. The minor of $M$ with row and column indices $0,1,2$ has a real zero in $(0,1)$ and therefore is not stable.
\end{theorem}
\begin{proof}
The matrices are real symmetric. For $v=(a,b,c)^T$, direct expansion gives
\begin{align*}
v^TD_2v&=13\left(a+\frac9{13}b+\frac7{13}c\right)^2
 +\frac{10}{13}\left(b+\frac15c\right)^2+\frac15c^2,\\
v^TD_3v&=5\left(a+\frac75b+\frac85c\right)^2
 +\frac65\left(b+\frac23c\right)^2+\frac23c^2.
\end{align*}
Each form is nonnegative, and equality forces successively $c=0$, $b=0$, and $a=0$. Thus both are strictly positive for nonzero vectors. Also $v^TD_1v=a^2+2b^2+c^2>0$ for $v\ne0$.

Put $P=\det(I+xD_1+yD_2+zD_3)$. Its complete expansion in $y,z$ is
\begin{align*}
P={}&2x^3+5x^2+4x+1+(49x^2+79x+30)z+(28x+22)z^2+4z^3\\
 &+y\bigl(41x^2+65x+24+(254x+164)z+64z^2\bigr)\\
 &+y^2(19x+16+62z)+2y^3.
\end{align*}
Consequently the selected submatrix is exactly
\[
 \begin{pmatrix}
2x^3+5x^2+4x+1&49x^2+79x+30&28x+22\\
41x^2+65x+24&254x+164&64\\
19x+16&62&0
\end{pmatrix}.
\]
Its determinant is
\[
\Delta(x)=-4(3076x^3+3012x^2-99x-440).
\]
We have $\Delta(0)=1760>0$ and $\Delta(1)=-22196<0$. Since $\Delta$ is continuous, the intermediate value theorem gives $x_0\in(0,1)$ with $\Delta(x_0)=0$. Viewed as a complex number, $x_0$ has positive real part. Hence the actual coefficient minor is not Hurwitz stable.
\end{proof}
The formal module \texttt{CoefficientMinor.lean} represents $P$ in the nested polynomial ring $\mathbb R[x][z][y]$, so its two coefficient operations are precisely $[z^j][y^i]$. It proves the complete determinant expansion before extracting coefficients. The natural value maps from \texttt{Fin 3} select the distinct increasing row and column indices $0,1,2$. The main theorem \texttt{CoefficientMinorCounterexample.source\_counterexample} combines all three strict positive-definiteness assertions with nonstability of the determinant of this selected submatrix. The complex nonvanishing definition is refuted by transporting the certified real root through the real-to-complex homomorphism.

Shivam Patel already posted this counterexample on MathDB. The matrices $D_2,D_3$ also occur in Fisk's example immediately following Conjecture 21.90; the posted solution replaces that example's semidefinite $D_1$ by the strictly positive definite matrix above. Fisk's nearby assertion that its displayed minors are stable is not used: the present certificate proves all computations independently. We claim verification of a full counterexample to the universal conjecture, not novelty or first formalization. The exact post and solution identifiers, snapshot hashes, and reproduction instructions accompany the module.
\begin{thebibliography}{9}
\bibitem{fisk} S. Fisk, \emph{Polynomials, roots, and interlacing}, arXiv:math/0612833, Chapter 21, definition at printed p.~619 and Conjecture 21.90 at printed p.~650, \url{https://arxiv.org/abs/math/0612833}.
\bibitem{patel} S. Patel, posted solution to MathDB 380287, \url{https://mathdb.com/p/380287}. Retrieved full post and identifier are supplied with the certificate.
\end{thebibliography}
\end{document}
